\newpage
\vspace{3cm}

\chapter*{Abstract}
Productive Failure (PF) is a learning design in which students attempt a challenging problem before they are taught the target concept. A good PF task lets students start from what they already know, invites several approaches, and lets familiar methods fall short so that the target concept becomes necessary. Such tasks are costly to design, and only a few published examples exist. While a large language model (LLM) could be prompted to write PF tasks in one pass, the success depends on whether the model knows what students at a given grade can and cannot yet do. In our replication, a rubric judge rated one-shot GPT-5 mini tasks at 14.6 of 15 points, yet familiar methods fell short on only $10\%$ of them.

To address this gap, we have developed \textsc{StructPF}, a three-stage pipeline that follows the decisions of a PF designer. We first collect the students' prior knowledge from curriculum data. Then, we plan which familiar method should fall short and verify the task numbers with calculations. Finally, we write the task text for students. Since the rubric could not separate working PF tasks from flawed ones, we also developed \textsc{PF Eval}, which tests the PF requirements on the finished task with five measures and a solver restricted to prior knowledge. On 29 reference tasks with known strengths and flaws, \textsc{PF Eval} matches every expected outcome.

We compared \textsc{StructPF} with one-shot generation in two studies, one with 358 tasks across four Grade~7 settings and one across 30 curriculum topics from Grades~5 to~10. Familiar methods fall short on $80\%$ of \textsc{StructPF} tasks compared with $10\%$ of one-shot tasks, and $57\%$ compared with $8.5\%$ pass all five measures. On 17 of the 30 topics, the one-shot prompt listed the target concept itself as prior knowledge, while curriculum grounding removed every such case. \textsc{StructPF} still trails the published expert tasks, of which 11 of 12 pass, and simulated students did not reproduce the approaches of real Grade~7 classrooms. These results suggest that PF generation is limited by knowledge about learners rather than by model size, and the generated tasks still need to be tested with real students.
